% ==========================================================
% titles/title-03.tex -- Title III: Governance Structure
% ==========================================================

\ConstTitle{Governance Structure}{title:governance-structure}

\Article{Local Democratic Governance}{art:local-democracy}

\ArtSection{Scope}{art:local-democracy:s-scope}
All decisions whose effects are substantially confined to a local community are made through direct local democracy.

\ArtSection{Direct Participation}{art:local-democracy:s-direct-participation}
All residents of a community participate directly in decisions affecting only that community. Decisions are made by a simple majority, subject to the Minority Impact Assessment described in \ref{art:local-democracy:s-minority-impact-assessment}.

\ArtSection{Facilitation and Rotation}{art:local-democracy:s-facilitation-rotation}
Facilitation roles at the local level rotate among residents on terms selected by lottery, not to exceed one year. Communities determine the specific term length within this maximum through their local democratic process. No person may serve consecutive facilitation terms without an intervening term served by another resident. There is no professional political class at the local level. All deliberations are fully public and recorded in real time by the \gls{gls:transparency-engine}.

\ArtSection{Minority Impact Assessment}{art:local-democracy:s-minority-impact-assessment}
Before any local decision takes effect, an independent assessment determines whether it significantly harms an identifiable minority. If so, the decision is suspended pending review by a regional Commons body. The assessment is conducted by the \gls{gls:minority-impact-assessment-body} and is fully transparent.

The \gls{gls:minority-impact-assessment-body} is a standard body of twenty-one members serving three-year terms, stratified to include meaningful representation from historically marginalized communities. It assesses distributional impacts of governance decisions before they take effect. Its assessments are advisory but mandatory in procedure -- no decision takes effect without a completed assessment and formal published response.

\ArtSection{Local Autonomy}{art:local-democracy:s-local-autonomy}
Communities have genuine autonomy within the foundational constraints of this Constitution. National governance may not override local decisions that do not violate foundational principles or impose harm on other communities.

\ArtSection{Community Justice Bodies}{art:local-democracy:s-community-justice-bodies}
\gls{gls:community-justice-bodies} are embedded within local democratic governance. Their composition, selection, terms, and mandate are established under Article \ref{art:justice-restoration}, \ref{art:justice-restoration:s-community-justice-bodies}.

\Article{The Sortition Legislature}{art:sortition-legislature}

\ArtSection{Composition and Selection}{art:sortition-legislature:s-composition-selection}
The \gls{gls:sortition-legislature} consists of six hundred persons selected by stratified random lottery as described in Article \ref{art:sortition-design}. Terms are three years, staggered so that two hundred members rotate out each year. The bootstrapping mechanism is described in Article \ref{art:sortition-design}, \ref{art:sortition-design:s-bootstrapping-first-legislature}. Legislators are subject to the conflict of interest and sector exclusion provisions applicable to all governance roles under Article \ref{art:governance-principles}, \ref{art:governance-principles:s-conflict-interest-sector}.

\ArtSection{Deliberative Process}{art:sortition-legislature:s-deliberative-process}
The Legislature follows a structured deliberative process for all significant legislation:
\begin{itemize}
  \item Expert briefings from multiple perspectives provided by the \gls{gls:expert-advisory-corps}
  \item Small-group deliberation in randomly assigned group
  \item Full-assembly reporting and cross-group deliberation
  \item Minority impact assessment for all legislation
  \item Final deliberation and decision
\end{itemize}

\ArtSection{The Expert Advisory Corps}{art:sortition-legislature:s-expert-advisory-corps}
A permanent \gls{gls:expert-advisory-corps} serves to advise the Legislature. It is a standard body of domain specialists serving three-year terms selected from the relevant \gls{gls:qualification-commons} pools. Members are explicitly forbidden from lobbying, advocacy, and post-service employment with entities that have interests in their advisory domain.

\ArtSection{Post-Service Employment}{art:sortition-legislature:s-post-service-employment}
A person is permanently prohibited from taking new employment with any entity that had active interests in legislative decisions the person made during their service, where the new role is connected to those decisions. All persons are prohibited from monetizing access to governance.

\ArtSection{Bad Faith and Removal}{art:sortition-legislature:s-bad-faith-removal}
Service requires genuine participation in good faith. Removal for bad faith requires authorization by the \gls{gls:constitutional-court} upon petition from the \gls{gls:citizen-oversight-panels} embedded within the \gls{gls:sortition-legislature} with a high evidentiary standard. The bad faith standard applies to process obstruction and corruption, never to substantive positions.

\ArtSection{Multilingualism}{art:sortition-legislature:s-multilingualism}
All legislative proceedings are conducted with real-time interpretation in all languages spoken by more than one percent of the Commonwealth population. Official documents and civic materials are available in all threshold languages.

\Article{The Constitutional Court}{art:const-court}

\ArtSection{Composition}{art:const-court:s-composition}
The \gls{gls:constitutional-court} consists of twenty-one justices selected by stratified sortition from a qualification pool maintained by the \gls{gls:qualification-commons} under Article \ref{art:sortition-design}, \ref{art:sortition-design:s-qualification-commons}, assessed through blind review of collected written analyses, not credential gatekeeping. Terms are seven years, staggered, and non-renewable.

\ArtSection{Function}{art:const-court:s-function}
The Court reviews legislation for compliance with this Constitution, resolves conflicts between governance layers, and interprets foundational principles in specific disputes. Constitutional override requires a supermajority of fourteen of twenty-one justices. Decisions touching the three hard limits require convergence of a Court supermajority, a Legislative supermajority, \gls{gls:meta-review-chamber} approval, and popular referendum.

\ArtSection{The Meta-Review Chamber}{art:const-court:s-meta-review-chamber}
A permanent \gls{gls:meta-review-chamber}, a standard body of thirty members serving one-year terms, reviews all \gls{gls:constitutional-court} decisions for quality of reasoning and consistency. The Chamber publishes public assessments of each decision and may trigger re-hearing by supermajority vote. It cannot overturn decisions but may compel reconsideration.

\ArtSection{Public Constitutional Petition}{art:const-court:s-public-constitutional-petition}
If five percent of the full civic membership formally petition against a specific \gls{gls:constitutional-court} decision within one year of its issuance, a new sortition court and chamber is constituted with no overlap in membership from the original bodies to re-examine the decision de novo. This mechanism applies only to \gls{gls:constitutional-court} decisions, cannot be used to circumvent the three hard limits, cannot be triggered more than once for the same decision, and does not suspend the original decision during re-examination unless the new court so orders. Petitions are submitted and verified through the \gls{gls:civic-initiative-infrastructure} under Article \ref{art:civic-initiative-infrastructure}.

\Article{The Constitutional Assembly and Generational Reaffirmation}{art:const-assembly}

\ArtSection{The Generational Cycle}{art:const-assembly:s-generational-cycle}
Every twenty-five years, a \gls{gls:constitutional-assembly} of one thousand persons, selected by stratified sortition, convenes for a minimum of six months to review every foundational principle and either reaffirm, revise, or repeal it. The Assembly determines its duration within this range by supermajority vote within its first thirty days of convening, based on its assessment of the scope of review required. No generation is permanently bound by the decisions of its predecessors.

\ArtSection{Thresholds for Change}{art:const-assembly:s-thresholds-change}
Changes to foundational principles require eighty percent approval within the Assembly, followed by ratification by sixty percent approval in a national direct vote.

\ArtSection{The Three Hard Limits}{art:hardlimits}
The Anti-Domination Principle, the Material Sufficiency Floor, and the Anti-Entrenchment Principle may not be abolished or diminished by any \gls{gls:constitutional-assembly}. They may be expanded. The process itself may be changed only by ninety percent Assembly approval plus national ratification.

\ArtSection{Extraordinary Assembly}{art:const-assembly:s-extraordinary-assembly}
If the \gls{gls:constitutional-court} (Article \ref{art:const-court}), the \gls{gls:citizen-oversight-panels} (Article \ref{art:citizen-oversight-panels}), and the \gls{gls:transparency-engine} (Article \ref{art:transparency-engine}) independently identify systemic constitutional failure beyond normal correction, they may jointly trigger an Extraordinary \gls{gls:constitutional-assembly} outside the normal twenty-five-year cycle. Because the Citizen Oversight Panels exist as many independent bodies rather than one, they speak for the purpose of this Section only through a two-thirds vote of the \gls{gls:cross-panel-coordination-body}'s representatives formally attesting that this threshold has been met; this attestation power is limited to this Section and does not otherwise expand the Cross-Panel Coordination Body's authority under Article \ref{art:citizen-oversight-panels}, \ref{art:citizen-oversight-panels:s-cross-panel-coordination}. Because the Transparency Engine operates without human discretion, its signal for the purpose of this Section is Anomaly Detection System reports sustained above a defined threshold over a defined period, as set by the \gls{gls:sortition-legislature} and published in advance. The mechanism for formal signaling and joint convening is published by the \gls{gls:electoral-commons} and logged through the \gls{gls:transparency-engine}.